\documentclass{article}
\usepackage{amsmath}
\usepackage{chicago}
\usepackage{setspace}
\usepackage{geometry}

\setcounter{MaxMatrixCols}{10}
\geometry{margin=1in}

\begin{document}

\title{\textsf{Learning, Commitment and Monetary Policy:}\\
\textsf{\ \ the case for partial commitment}}
\author{\textsf{George A. Waters\thanks{%
gawater@ilstu.edu}} \\
\textsf{Department of Economics}\\
\textsf{Campus Box 4200}\\
\textsf{Illinois State University}\\
\textsf{Normal, IL 61761-4200}}
\date{\textsf{June 21, 2008}}
\maketitle

\begin{abstract}
This paper examines a class of interest rate rules, studied in Evans and
Honkapohja (2003, 2006), that respond to public expectations and to lagged
variables. \ Their work is extended by considering varying levels of
commitment that correspond to varying degrees of response to lagged output.
\ Under commitment the policymaker adjusts the nominal rate with lagged
output to impact public expectations. \ Within this class of rules, I
provide a condition for non-explosive and determinate solutions. \
Expecatational stability obtains for any non-negative response to lagged
output. \ Simulation results show that modified commitment, as advocated by
Blake (2001), is best under least squares learning. \ However, in the
presence of parameter uncertainty and/or measurement error in the
policymaker's data on public expectations, the best policy is one of partial
commitment, where the response to lagged output is less than under modified
commitment. \ The case for partial commitment is strengthened if the gain
parameter in the learning mechanism is high, which can be interpreted as the
use of few lags by public agents in the formation of expectations or as an
indication of low credibility of the policymaker. \ The appointment of a
conservative central banker ameliorates these concerns about modified
commitment.

\textbf{Keywords: \ }Learning, Monetary Policy, Interest Rate Rules,
Determinacy, Expectational Stability

\textbf{JEL classification:} \ E52, E31, D84
\end{abstract}

\pagebreak

\begin{doublespace}

\section{Introduction}

The gains from commitment for a monetary policymaker have been established
both theoretically (Barro and Gordon (1983) and Clarida, Gali and Gertler
(1999)) and empirically (McCallum and Nelson (2004))\footnote{%
Time consistency of monetary policy has been a relevant issue going back to
Kydland and Prescott (1977) and continuing with Woodford's (1999b, 2003)
discussion of a \textit{timeless perspective.}}. \ Another strand of work on
monetary policy (Howitt (1992) and Bullard and Mitra (2001))\footnote{%
Other examples include Honkapohja and Mitra (2003) and Beechey (2004).}
advocates the view that monetary policy rules should be learnable, meaning
they lead to dynamically stable systems when public expectations deviate
from full rationality. \ Setting policy under commitment implies that the
policymaker acts to impact public expectations, so the way the public forms
expectations could be critical to the outcome of such an approach. \ The
present work addresses the question of whether there are gains from
commitment under learning.

Interest rate rules that respond explicitly to public expectations, studied
by Evans and Honkapohja (2003, 2006), have the desirable properties of
determinacy and expectational stability under both commitment and
discretionary policies\footnote{%
Evans and McGough (2005) study determinancy and expectational stability for
different policy rules and specifications of the Phillip's curve.}. \ Within
this class, the optimal policy rule under rational expectations is identical
in both cases with one exception. \ Under commitment, policy responds to
lagged values of output while under discretion it does not. \ As Woodford
(1999a) emphasizes, if public agents are forward looking, commitment
requires that policy be \textit{history dependent}. \ The more the
policymaker wants to impact public expectations, the greater the response to
lagged variables.

Since the existence of gains to commitment under learning is an open
question, I study a continuum of responses to lagged output that includes
the optimal rational expectations discretionary and commitment policies as
special cases. \ There are other reasons to consider such a range of
policies. \ First, if the policymaker has biased estimates of the model
parameters, he or she may over- or under-react to lagged output under
commitment. \ I also consider the possibility that the policymaker's
knowledge of public expectations is imprecise.\ \ For the expanded set of
rules, I provide a condition under which there are non-explosive and
determinate rational expectations equilibria. \ These equilibria are
expectationally stable for any non-negative response to lagged output.

Simulations of the model under learning allow a comparison of policy
outcomes for various specifications of policy and the learning mechanism. \
Public expectations are formed with recursive least squares using constant
gain. \ The gain parameter is a key component, indicating agents'
sensitivity to new information, and has natural interpretations in terms of
the number of lags used in the estimation, the rationality of the agents and
the credibility of the policymaker.

One can regard any policy response to lagged output as a type of commitment
in that it is an attempt to influence expectations. \ The commitment optimum
under rational expectations, as derived\footnote{%
Woodford (1999a) and Blake (2001) refer to this as the timeless perspective
policy.} in Evans and Honkapohja (2006), is referred to as \textit{full
commitment}. \ Blake (2002) argues for a particular\footnote{%
Jensen and McCallum (2002) also discuss this point.} intermediate response
to lagged output called \textit{modified commitment }for the present study.
\ The results from the simulations demonstrate that, under least squares
learning, the policy outcome under modified commitment minimizes the
policymaker's loss in comparison to full commitment, discretion and all
other settings for the response to lagged output. \ However, the results
also indicate that if there is uncertainty about the parameters of the
model, the policymaker should not respond to lagged output as strongly as
the full or modified commitment settings suggest. \ I refer to such a policy
prescription as \textit{partial commitment}.

Implementing expectations based interest rate rules raises questions about
the policymaker's knowledge of public expectations. \ Evans and Honkapohja
(2006) note that the introduction of white noise errors to the policymaker's
estimates of public expectations does not alter their results concerning
expectational stability. \ However, whether such errors affect the policy
outcomes under learning remains an open question.

The results presented here show that the introduction of such errors
substantiate the case for partial commitment. \ In the presence of errors in
the policy rule, the response to lagged output should be weaker than the
full or modified commitment recommendations. \ Moreover, for higher levels
of the gain parameter, the optimal response to lagged output falls further.
\ If agents are placing heavy emphasis on recent information, possibly due
to a lack of credibility of the policymaker, partial commitment is the best
approach under learning.

The results are verified for different parameter values found in the
literature. \ The model studied here is the standard New Keynesian approach
where the policymaker faces a trade-off between stabilizing output and
inflation, described in Clarida, Gali and Gertler (1999). \ Results are very
similar for different values of the model parameters, but they are sensitive
to changes in the policymaker's preference parameter measuring his or her
relative interest in these two goals. \ 

A significant factor underlying the optimality of partial commitment is the
asymmetry of the losses across the magnitude of the response to lagged
output, in that an excessive response can leads to very poor outcomes from
the policymaker's perspective. \ This paper provides support for a
conservative central banker (Rogoff 1985) since a lower emphasis on output
stabilization diminishes the impact of this asymmetry and increases the
range for non-explosive and determinate policies.

The learning rule in the simulated model uses constant gain, meaning the
emphasis agents place on new information does not change over time. \
Therefore, learning has a significant impact on the dynamics of the model at
all times. \ Orphanides and Williams (2002, 2006) study an alternative model
of monetary policy under constant gain learning and show that the nature of
optimal policy changes when expectations deviate from full rationality. \
Constant gain\footnote{%
Honkapohja and Mitra (2002) give theoretical stability results under bonded
memory learning. \ Carceles-Poveda and Giannitsarou (2006) use constant gain
to study asset price dynamics.} is an important aspect of the dynamics in
Bullard and Cho (2005), who show how the economy can shift to a deflationary
equilibrium when agents continually update their estimate of the
policymaker's inflation target. \ 

The approach to expectation formation employed here follows a large
literature on learning and convergence to rational expectations that relaxes
the extreme informational requirements on agents required by rational
expectations. \ Bray (1982) is one of the earliest efforts, while Marcet and
Sargent (1989) give convergence criteria for least squares learning. \ Evans
and Honkapohja (2001) provide an overview of this literature and define
expectational stability for a wide variety of learning mechanisms and
models. \ Howitt (1992) contains one of the earliest suggestions that
monetary policy rules should be stable under learning. \ More recently,
Bullard and Mitra (2002) and Honkapohja and Mitra (2004) have studied the
expectational stability of various interest rate rules for monetary policy.

The paper is organized as follows. \ Section 2 describes the model, and
derives conditions for optimal policy for varying levels of commitment. \
Section 3 shows the computation of the rational expectations equilibrium,
and section 4 introduces the expectations based interest rate rule and
demonstrates the condition for non-explosiveness and determinacy. \ Section
5 describes the learning mechanism and the proposition concerning
expectational stability. \ Section 6 reports the initial simulation results
while Section 7 discusses those for the model with errors in the policy rule
and examines the simulation results for alternative parameters values. \
Section 9 concludes.

\section{The Model}

The following New Keynesian model has become standard for the study of
interest rate rules in monetary policy. \ It has micro foundations described
in Woodford (2003) including price stickiness that allows the policymaker to
play a stabilizing role in the economy.%
\begin{equation}
x_{t}=-\varphi \left( i_{t}-E_{t}^{\ast }\pi _{t+1}\right) +E_{t}^{\ast
}x_{t+1}+g_{t}  \label{IS}
\end{equation}%
\begin{equation}
\pi _{t}=\lambda x_{t}+\beta E_{t}^{\ast }\pi _{t+1}+u_{t}  \label{Phillip's}
\end{equation}%
These equations are expectations augmented IS and Phillip's curve
relationships. \ The variables $x_{t}$ and $\pi _{t}$ are the deviations of
output and inflation from their target values. \ The notation $E_{t}^{\ast }$
indicates private sector expectations formed in time $t$ where the $\left(
\ast \right) $ is used to show that expectations might not be rational. \
The policymaker controls the nominal interest rate $i_{t}$. \ The parameters 
$\varphi ,\lambda ,\beta $ are all positive and the discount rate $\beta $
is such that $\beta <1$. \ The stochastic terms $g_{t}$ and $u_{t}$ both
have autoregressive structure.%
\begin{eqnarray}
\left( 
\begin{array}{c}
g_{t} \\ 
u_{t}%
\end{array}%
\right) &=&F\left( 
\begin{array}{c}
\widehat{g}_{t-1} \\ 
u_{t-1}%
\end{array}%
\right) +\left( 
\begin{array}{c}
\widetilde{g}_{t}+\widetilde{w}_{t} \\ 
\widetilde{u}_{t}%
\end{array}%
\right)  \label{shocks} \\
\text{for \ \ }F &=&\left( 
\begin{array}{cc}
\mu & 0 \\ 
0 & \rho%
\end{array}%
\right)  \notag
\end{eqnarray}%
The parameters $\mu $ and $\rho $ lie in the interval (-1, 1), and the
shocks $\widetilde{u}_{t},$ $\widetilde{w}_{t}$ and $\widetilde{g}_{t}$ and
are \textit{iid }with standard deviations $\sigma _{u},$ $\sigma _{w}$ and $%
\sigma _{g}$, respectively. \ The structure of $g_{t}$ follows McCallum and
Nelson (2004) who decompose this term into a preference shock $\widetilde{g}%
_{t}$ and an AR(1) process $\widehat{g}_{t}=\mu \widehat{g}_{t-1}+\widetilde{%
w}_{t}$ that accounts for uncertainty about the evolution of the natural
rate that enters into the forecast error for output in the IS equation (\ref%
{IS}) and represents the portion of the demand shocks about which the public
has some information to use for forecasting.

The policymaker sets the nominal interest rate to stabilize the endogenous
variables. \ Formally the task is to set $i_{t}$ to minimize the loss
function%
\begin{equation}
L=E_{t}\dsum\limits_{s=0}^{\infty }\beta ^{s}\left( \pi _{t+s}^{2}+\alpha
x_{t+s}^{2}\right) ,  \label{Loss}
\end{equation}%
assuming rational expectations for the following derivation.\ \ The
parameter $\alpha $ measures the relative importance of output and inflation
stabilization for the policymaker, the value $\alpha =0$ corresponding to
pure inflation targeting. \ Minimizing the loss function\footnote{%
The present approach follows Evans and Honkapohja (2004), Woodford (1999b)
and Clarida, Gali and Gertler (1999).} over $x_{t+s}$ and $\pi _{t+s}$ under
the constraint of the Phillip's curve (\ref{Phillip's}) yields the following
first order conditions.%
\begin{equation}
E_{t}\left( 2\alpha x_{t+s}+\lambda \omega _{t+s}\right) =0\text{ \ \ \ for }%
s=0,1,2,3...  \label{FOC 1}
\end{equation}%
\begin{equation}
E_{t}\left( 2\pi _{t+s}+\omega _{t+s-1}-\omega _{t+s}\right) =0\text{\ \ \
for }s=1,2,3...  \label{FOC 2}
\end{equation}%
\begin{equation}
2\pi _{t}-\omega _{t}=0  \label{FOC 3}
\end{equation}%
\qquad The variable $\omega _{t+s}$ is the Lagrange multiplier on the
constraint for each $s=0,1,2,3....$ \ Optimal policy in time $t$ is governed
by (\ref{FOC 3}) but policy in succeeding periods is determined by the above
condition (\ref{FOC 2}). \ The time consistency problem is evident since,
when the policymaker solves the problem in the period $t+1$, the policy
given by (\ref{FOC 3}) will be different \ than the policy prescribed by (%
\ref{FOC 2}) in the period $t$ solution.

To find the condition for policy under discretion, the policymaker adopts
the optimal policy each period, solving out $\omega _{t}$ for $s=0$ from the
first order conditions (\ref{FOC 1}) and (\ref{FOC 3}).%
\begin{equation}
\lambda \pi _{t}+\alpha x_{t}=0  \label{Discretion cond}
\end{equation}%
Under discretion, the policymaker takes expectations to be fixed and ignores
(\ref{FOC 2}), see Clarida, Gali and Gertler (1999) for a detailed
discussion. \ However, if the policymaker has a "timeless perspective"
(Woodford 1999b) he or she uses the policy that would have been prescribed
in past periods, ignoring condition (\ref{FOC 3}), using (\ref{FOC 1}) and (%
\ref{FOC 2}) for $s=0$. \ In this case, which I refer to as \textit{full
commitment}, the policymaker acts to influence future expectations yielding
the condition%
\begin{equation}
\lambda \pi _{t}+\alpha \left( x_{t}-x_{t-1}\right) =0.
\label{Commitment cond}
\end{equation}

Policy under commitment differs from discretion in its response to the
previous period's output gap. \ The reason can be seen in (\ref{FOC 2}),
where the lagged term appears because the policymaker reacts to the expected
inflation in the Phillip's curve (\ref{Phillip's}). \ Under discretion,
however, the policymaker ignores the change in expected inflation, and
lagged output does not enter the policy condition (\ref{Discretion cond}). \ 

There are several reasons to question full commitment, meaning the
policymaker should not be restricted to these extreme cases. \ One is
potential bias in the policymaker's estimates of model parameters, while
another is the possibility that the policymaker may not have full
information about public expectations. \ Also, gains to commitment depend on
a forward looking public, so if the public is using a learning mechanism to
form expectations, the policymaker may not want to fully commit. \ Waters
(2005) shows that discretion is superior to commitment under adaptive
expectations, where the public forecast of a variable is simply a weighted
average of the previous period's forecast and realized value. \ While the
least squares mechanism studied in the present paper is a more sophisticated
method of forming expectations, it remains an open question whether full
commitment is optimal. \ 

Furthermore, Blake (2001) argues that the first order condition (\ref{FOC 2}%
) improperly assumes certainty equivalence and he derives optimal policy
that differs from full commitment. \ Using a version of the loss function (%
\ref{Loss}) without discounting over a finite number of periods, he gives an
intuitive argument\footnote{%
He also checks the result with a formal derivation. Jensen and McCallum
(2002) also discuss this issue. \ Evans and McGough (2006) refer to this
approach as the MJB-alternative in their study of the New Keynsian model
with inertia} for the policymaker to respond to lagged output, but not to
the degree indicated by the full commitment condition (\ref{Commitment cond}%
). \ This approach captures the spirit of the timeless perspective, since
including discounting over-emphasizes the loss in initial periods. \ For
these reasons, we consider the general condition%
\begin{equation}
\lambda \pi _{t}+\alpha \left( x_{t}-\kappa x_{t-1}\right) =0.
\label{general cond}
\end{equation}%
The condition under discretion is a special case of (\ref{general cond})
where $\kappa =0,$ and full commitment is equivalent to setting $\kappa =1$.
\ 

The effect of Blake's (2001) approach without discounting is that the second
term in (\ref{FOC 2}) is now $\beta \omega _{t+s-1}$. \ In this case,
optimal policy corresponds to the condition (\ref{general cond}), where $%
\kappa $ equals the discount factor $\beta <1$, which I call \textit{%
modified commitment}. \ Woodford (2001) argues in favor of full commitment
by using a slightly different criterion than the loss function (\ref{Loss})
for evaluating policy. \ Under rational expectations the optimal level of $%
\kappa $ is either full commitment at $\kappa =1$ or modified commitment $%
\kappa =\beta $ depending on the approach to the derivation. \ I refer to
any setting for $\kappa $ such that $0<\kappa <\beta ,$ when the
policymaker's response to lagged inflation is less than for modified
commitment, as \textit{partial commitment}.

\section{RE Equilibria}

One can now solve for the unique minimum state variables\footnote{%
See McCallum (1983, 1997).} rational expectations equilibrium (REE) for a
given $\kappa $ in the general condition (\ref{general cond}). \ Postulating
solutions of the form%
\begin{eqnarray}
x_{t} &=&b_{x}x_{t-1}+c_{x}u_{t}  \label{RE sols} \\
\pi _{t} &=&b_{\pi }x_{t-1}+c_{\pi }u_{t}  \notag
\end{eqnarray}%
implies that $E_{t}\pi _{t+1}=b_{\pi }\left( b_{x}x_{t-1}+c_{x}u_{t}\right)
+c_{\pi }\rho u_{t}$. \ Using the method of undetermined coefficients with (%
\ref{general cond}) and (\ref{Phillip's}) yields%
\begin{eqnarray}
b_{\pi } &=&\lambda b_{x}+\beta b_{\pi }b_{x}  \label{coeff sols} \\
\lambda b_{\pi } &=&\alpha (\kappa -b_{x})  \notag \\
c_{\pi } &=&\lambda c_{x}+\beta \left( b_{\pi }c_{x}+c_{\pi }\rho \right) +1
\notag \\
\lambda c_{\pi } &=&-\alpha c_{x}.  \notag
\end{eqnarray}%
Combining the first two equations from (\ref{coeff sols}) gives us%
\begin{equation}
\beta b_{x}^{2}-\left( \frac{\lambda ^{2}}{\alpha }+\beta \kappa +1\right)
b_{x}+\kappa =0  \label{quad cond}
\end{equation}%
The larger root of the quadratic equation (\ref{quad cond}) is always such
that $b_{x}>1$ so the smaller root is the potentially non-explosive solution
for $b_{x}$ in (\ref{RE sols}) and is the focus throughout the paper. \
Given $b_{x}$, the solution for $b_{\pi }$ is determined by the first and
third equations in (\ref{coeff sols}). \ Note that the discretionary value $%
\kappa =0$ admits the solution $b_{x},b_{\pi }=0,$ which corresponds\ to the
fact that the minimum state variables solution in this case depends only on
the supply shock $u_{t}$, as in Evans and Honkapohja (2003). Further, for $%
\kappa >0,$ the relevant solutions for $b_{x}$ and $b_{\pi }$ are real and
positive (see Appendix B). \ Hence, public expectations respond to lagged
output only when the policymaker does, demonstrating the policymaker's
influence on expectations under commitment. \ The solutions for the
coefficients under full commitment for $\kappa =1$ correspond to those in
Evans and Honkapohja (2006). \ Similarly, the third and fourth equations
from (\ref{coeff sols})\ may be combined to give%
\begin{equation*}
c_{x}=\lambda \left[ \alpha \beta \rho -\alpha -\lambda ^{2}-\beta \lambda
b_{\pi }\right] ^{-1}
\end{equation*}%
and then the fourth equation in (\ref{coeff sols}) determines $c_{\pi }$,
completing the solution.

\section{Expectations Based Interest Rate Rules}

Evans and Honkapohja (2003, 2006) advocate for monetary policy to be
conducted with interest rate rules that respond explicitly to public
expectations. \ One can compute the optimal form under rational expectations
for the expectations based interest rate rule associated with the policy
condition (\ref{general cond}). \ This section also provides a condition for
the non-explosive and determinate solutions, while the following section on
learning addresses expectational stability.

Using (\ref{general cond}) to substitute for inflation in the Phillip's
curve (\ref{Phillip's}) gives%
\begin{equation*}
x_{t}=\lambda \left( \lambda ^{2}+\alpha \right) ^{-1}\left( \alpha \kappa 
\lambda ^{-1}x_{t-1}-\beta E_{t}\pi _{t+1}-u_{t}\right) .
\end{equation*}%
Substituting out $x_{t}$ in the IS equation with the above expression yields
the interest rate rule%
\begin{equation}
i_{t}=\delta _{L}x_{t-1}+\delta _{\pi }E_{t}^{\ast }\pi _{t+1}+\delta 
_{x}E_{t}^{\ast }x_{t+1}+\delta _{g}g_{t}+\delta _{u}u_{t}
\label{policy rule}
\end{equation}%
where the parameters are%
\begin{eqnarray*}
\delta _{L} &=&-\varphi ^{-1}\left( \lambda ^{2}+\alpha \right) ^{-1}\alpha
\kappa ,\text{ \ \ }\delta _{\pi }=1+\varphi ^{-1}\left( \lambda ^{2}+\alpha
\right) ^{-1}\lambda \beta \\
\delta _{x} &=&\varphi ^{-1},\text{ \ }\delta _{g}=\varphi ^{-1},\text{ \ \ }%
\delta _{u}=\varphi ^{-1}\left( \lambda ^{2}+\alpha \right) ^{-1}\lambda .
\end{eqnarray*}%
Note that the extent to which the interest rate responds to lagged output,
shown by $\delta _{L}$, depends on $\kappa $, but the other parameters in
the rule do not. Under discretion, $\kappa $ is zero and $i_{t}$ is
unaffected by $x_{t-1}$, but for any other value, including the full,
modified and partial commitment settings, the interest rate responds
directly to $x_{t-1}$. \ 

The interest rate rule (\ref{policy rule}) demonstrates the potential impact
of parameter uncertainty on the conduct of policy. \ For example, if the
policymaker's estimate $\varphi _{p}$ is less than the actual $\varphi $,
the magnitude of the response of $i_{t}$ to $x_{t-1}$ will be too large,
which is equivalent to setting $\kappa $ higher than intended, and
simulation results indicate that such a policy could lead to a large loss%
\footnote{%
If the policymaker has a biased estimate of $\lambda $, the change in policy
would be equivalent to the change caused by varying the policymaker
preference parameter $\alpha $.}. \ For this reason, it is important to
study the behavior of the model when $\kappa >1$, even though the
policymaker would never knowingly make such a choice. \ To guard against
such a possibility, the policymaker may not fully commit and may choose a
smaller response to lagged output than would be indicated by (\ref%
{Commitment cond}). \ 

The next tasks are to determine conditions under which interest rate rules
of the form (\ref{policy rule}) are non-explosive, determinate and
expectationally stable. \ A determinate equilibrium is unique in that there
are no other similar equilibria based on extraneous variables. \
Expectational stability implies that the expectations of agents using a
reasonable learning rule converge to the REE. \ All of these features of
interest rate rules are desirable for a policymaker trying to stabilize
endogenous variables.

For analysis of determinacy, expectational stability and for simulations, it
is convenient to write the New Keynesian model with (\ref{IS}), (\ref%
{Phillip's}) and (\ref{policy rule}) in matrix form.%
\begin{eqnarray}
\left( 
\begin{array}{c}
x_{t} \\ 
\pi _{t}%
\end{array}%
\right) &=&\left( 
\begin{array}{cc}
\left( \lambda ^{2}+\alpha \right) ^{-1}\alpha \kappa & 0 \\ 
\left( \lambda ^{2}+\alpha \right) ^{-1}\lambda \alpha \kappa & 0%
\end{array}%
\right) \left( 
\begin{array}{c}
x_{t-1} \\ 
\pi _{t-1}%
\end{array}%
\right)  \label{matrix model} \\
&&+\left( 
\begin{array}{cc}
0 & -\left( \lambda ^{2}+\alpha \right) ^{-1}\beta \lambda \\ 
0 & \left( \lambda ^{2}+\alpha \right) ^{-1}\beta \alpha%
\end{array}%
\right) \left( 
\begin{array}{c}
E_{t}^{\ast }x_{t+1} \\ 
E_{t}^{\ast }\pi _{t+1}%
\end{array}%
\right)  \notag \\
&&\text{ \ \ \ \ \ \ \ }+\left( 
\begin{array}{cc}
0 & -\left( \lambda ^{2}+\alpha \right) ^{-1}\lambda \\ 
0 & \left( \lambda ^{2}+\alpha \right) ^{-1}\alpha%
\end{array}%
\right) \left( 
\begin{array}{c}
g_{t} \\ 
u_{t}%
\end{array}%
\right)  \notag
\end{eqnarray}%
From the above, one can see that the policymaker acts so that the shock to
the IS equation $g_{t}$ does not affect output and inflation, but the supply
shock $u_{t}$ does. \ Note that if $\alpha =0$, meaning the policymaker
targets only inflation, the supply shock does not affect inflation. \ These
observations reflect the trade-off faced by the policymaker between output
and inflation stabilization that arises from the presence of the supply
shock.

The next proposition gives a single condition on $\kappa $ guaranteeing
non-explosivesness and determinacy. \ 

\begin{proposition}
Under rational expectations, for $\kappa $ such that $0\leq \kappa <1+\dfrac{%
\lambda ^{2}}{\alpha \left( 1-\beta \right) }$ ,

\begin{itemize}
\item the model defined by the policy condition (\ref{general cond}) and the
Phillip's curve (\ref{Phillip's}) has a non-explosive solution such that $%
b_{x}<1,$ and

\item the model given by (\ref{matrix model}) is determinate.
\end{itemize}
\end{proposition}

\begin{proof}
See appendix A.
\end{proof}

The condition for non-explosivenes applies to the matrix form (\ref{matrix
model}) as the policy rule (\ref{policy rule}) is derived from the general
condition (\ref{general cond}). \ The discretionary and full commitment
cases of $\kappa =0$ and 1 yield a determinate model, as shown in Evans and
 Honkapohja (2003, 2006), as do any intermediate values of $\kappa $
including modified and partial commitment, as well as some values of $\kappa
>1$. \ In the absence of extraneous variables in the simulated model,
indeterminacy should not be a crucial factor in the policy outcomes for the
simulations reported here, but given the number of candidate variables
(stock market indices, real estate market conditions) agents might consider
for use in forming expectations in practice, determinacy remains a desirable
property for a policy rule.

One might argue that values of $\kappa $ such that $\kappa >1$ are not
interesting as there is no justification for such policy. \ However, as
noted above, the policymaker with mistaken estimates of the model parameters
could easily overreact to lagged output. \ This study reports simulation
results using parameters from three papers in the literature, McCallum and
Nelson (2004), Clarida, Gali and Gertler (2000) and Woodford (1999),
referred to as MN04, CGG00 and W99, respectively. \ The following table
gives their estimates for the parameters in equations (\ref{IS}) and (\ref%
{Phillip's}).%
\begin{equation}
\begin{tabular}{|c|c|c|c|}
\hline
& $\beta $ & $\varphi $ & $\lambda $ \\ \hline
$\text{MN04}$ & $0.99$ & $0.4$ & $0.05$ \\ \hline
$\text{CGG00}$ & $0.99$ & $4.0$ & $0.075$ \\ \hline
$\text{W99}$ & $0.99$ & $\left( 0.157\right) ^{-1}$ & $0.024$ \\ \hline
\end{tabular}
\label{table}
\end{equation}%
There is quite a difference of opinion regarding $\varphi $ and $\lambda $
that could affect the response of the interest rate to lagged output,
represented by $\delta _{L}$ in (\ref{policy rule}). \ If the policymaker
fully commits ($\kappa =1.0$), places relatively low emphasis on output
stabilization with $\alpha =0.1$, and takes the parameter values from MN04
to be correct, the value for $\delta _{L}$ is $\delta _{L,MN}=-2.439$, but
if the true values are those in CGG00, this parameter should be $\delta 
_{L,CGG}\approx -0.237.$ \ Here, the policymaker would mistakenly set policy
corresponding to a value of $\kappa \approx 10.3$, which exceeds the bound
of $5.625$ for determinacy and non-explosiveness given in Proposition 1. \
So parameter uncertainty could lead the policymaker to set policy in an
explosive and indeterminate range. \ Evans and McGough (2006) make a similar
argument in the context of Taylor rules in the New Keynesian model.

Proposition 1 provides an additional rationale for the appointment of a
conservative central banker meaning one with a low weight on output
stabilization $\alpha $. \ Rogoff (1985) shows a conservative central banker
can diminish the problem of inflation bias\footnote{%
See Walsh (2003) for a full discussion of inflation bias.} in models where
the policymaker wants to raise output above the natural rate. \ In the
present framework, the policymaker cares only about deviations from the
natural rate, but a low $\alpha $ may still be desirable since it increases
the bound for determinacy and non-explosiveness, making such issues less
problematic.

\section{Learning}

The next task is to specify the learning mechanism agents use to form
expectations. \ As is common in the learning literature, I use recursive
least squares\footnote{%
As noted Orphanides and Williams (2002) is closely related. \ Fuhrer and
Hooker (1993) is another example of a monetary policy analysis using least
squares learning.}. \ This mechanism has a natural relationship with
expectational stability and can be used to simulate the model, which also
allows for a comparison of policy outcomes for varying parameter values. \
Further goals are to find a condition for expectational stability and to
give interpretations for the gain, a key parameter in the learning mechanism.

Assume agents update expectations using a model, called the perceived law of
motion (PLM), that corresponds in structure to the minimum state variables
solutions (\ref{RE sols}). \ Let the vectors $y_{t}$ and $v_{t}$\ be such
that $y_{t}=\left( x_{t},\pi _{t}\right) ^{\prime }$ and $v_{t}=\left(
g_{t},u_{t}\right) ^{\prime }$. \ Now, the PLM can be written as%
\begin{equation}
y_{t}=a+by_{t-1}+cv_{t}\text{.}  \label{PLM}
\end{equation}%
The REE corresponds to the case where $a=\left( 
\begin{array}{c}
0 \\ 
0%
\end{array}%
\right) $, $b=\left( 
\begin{array}{cc}
b_{x} & 0 \\ 
b_{\pi } & 0%
\end{array}%
\right) $ and $c=\left( 
\begin{array}{cc}
0 & c_{x} \\ 
0 & c_{\pi }%
\end{array}%
\right) $, but, under learning, agents do not know these values and must
update their estimate of $\xi _{t}=\left( a_{t},b_{t},c_{t}\right) $ over
time. \ The PLM (\ref{PLM}) with time subscripts for $a_{t},b_{t}$ and $%
c_{t} $ denotes agents beliefs at time $t$\ about parameter values and
implies that $E_{t}^{\ast }y_{t+1}=a_{t}+b_{t}E_{t}^{\ast
}y_{t}+c_{t}E_{t}^{\ast }v_{t+1}$ or%
\begin{equation}
E_{t}^{\ast }y_{t+1}=a_{t}+b_{t}\left( a_{t}+b_{t}y_{t-1}+c_{t}v_{t}\right)
+c_{t}Fv_{t}.  \label{exp}
\end{equation}%
This equation shows how agents use past information about the endogenous
variables, estimates about parameter values in the PLM and current values of
the shocks to form expectations.

Given such a description about expectations formation and the model (\ref%
{matrix model}), one can describe the actual law of motion (ALM) governing
the endogenous variables. \ Agents update the parameters to find $\xi _{t}$
and use the new shocks $v_{t}$ to form expectations $E_{t}^{\ast }y_{t+1}$
with (\ref{exp}), which then determines the contemporaneous value of $y_{t}$
via the model\footnote{%
Note that exogenous variables and parameter estimates dated time $t$ are
used to form expectations of future variables, but endogenous variables
dated time $t$ are not.} (\ref{matrix model}). \ One key question is whether
the parameters $\xi _{t}$ will converge over time to their values in the
REE, given by (\ref{coeff sols}). \ 

Evans and Honkapohja (2001) define expectational stability, which implies
the convergence of many learning mechanisms including least squares learning
(see appendix B). \ For the present model, one can show the following.

\begin{proposition}
The REE given by (\ref{RE sols}, \ref{coeff sols}) for model (\ref{matrix
model}) under the PLM (\ref{PLM}) is expectationally stable for any $\kappa
\geq 0$.
\end{proposition}

\begin{proof}
See appendix B.
\end{proof}

Expectational stability obtains for any level of response to lagged output,
providing further support for interest rate rules of the form (\ref{policy
rule}) and extending the expectational stability results of Evans and
Honkapohja (2003, 2006). \ The condition in Proposition 1 thereby guarantees
non-explosiveness, determinacy and expectational stability.

To formally define the least squares learning mechanism, let $z_{t}$ be such
that%
\begin{equation*}
z_{t}=\left( 1,y_{t-1}^{\prime },v_{t}^{\prime }\right) .
\end{equation*}%
The recursive least squares algorithm\footnote{%
See Evans and Honkapohja (2001) and Marcet and Sargent (1989) for details.}
is given by%
\begin{eqnarray}
R_{t} &=&R_{t-1}+\tau _{t}\left( z_{t-1}z_{t-1}^{\prime }-R_{t-1}\right)
\label{RLS} \\
\xi _{t} &=&\xi _{t-1}+\tau _{t}R_{t}z_{t-1}\left( y_{t-1}-\xi 
_{t-1}z_{t-1}\right) ^{\prime },  \notag
\end{eqnarray}%
where $R_{t}$ acts as an updated covariance matrix. \ The realization for $%
\xi _{t}$ then determines the expectations in (\ref{matrix model}), which
then determine $x_{t}$ and $\pi _{t}$. \ So the dynamics of the model are
defined by the recursive least squares updating equation (\ref{RLS}), the
expectations (\ref{exp}) derived from the PLM, the model (\ref{matrix model}%
) and the stochastic shocks.

A key parameter in the learning mechanism is the gain parameter $\tau _{t}$,
which indicates the emphasis that agents place on recent information. \ To
form an expectation, agents update their forecasting model (the PLM) from
the previous period using the new realizations of the endogenous variables.
\ A high gain indicates that agents' estimates of $\left(
a_{t,}b_{t},c_{t}\right) $ are more sensitive to recent information.

If the gain parameter is decreasing over time such that $\tau _{t}=1/t$,
then the updating equations (\ref{RLS}) are equivalent to recursive least
squares using all lags, given an appropriate initial condition. \ The
expectational stability result in Proposition 2 implies that recursive least
squares will converge to the REE from (\ref{RE sols}) and (\ref{coeff sols}%
), given that the PLM (\ref{PLM}) is of the same form as the minimum state
variables solutions (\ref{RE sols}) (see Evans and Honkapohja (2001b),
section 2.6 and chapter 10).

To study the impact of learning on policy, decreasing gain is not an
appropriate assumption, since the impact of learning on the dynamics is
decreasing over time. \ Following much of the learning literature, I adopt a
constant gain $\tau _{t}=\bar{\tau}$, the value of which has a number of
interrelated interpretations. \ Orphanides and Williams (2002) explain that
the learning mechanism (\ref{RLS}) with constant gain updates expectations
using a rolling window of past data. \ A lower gain parameter means that
agents are using more lags to make their estimates\footnote{%
A constant gain parameter $\bar{\tau}$\ corresponds to $2/\bar{\tau}$ \ lags
of data used.}. \ Furthermore, since least squares learning with decreasing
gain approaches the rational expectation equilibrium, one can think of the
gain parameter as a measure of the rationality of agents, lower gain
indicating a higher degree of rationality. \ A third interpretation can be
found in Waters (2007), who argues that low gain implies high credibility of
the policymaker. \ Given that the policymaker is trying to keep output and
inflation near certain targets, if agents have faith in the abilities of the
policymaker, they will not be excessively swayed by recent data and will
tend to keep their beliefs about the target values that they have formed
over time. \ Conversely, if agents do not trust the policymaker to keep the
economy near the targets or even to keep fixed targets, they will place
greater emphasis on recent realizations when forming expectations,
corresponding to high gain.

\section{Simulations}

Simulating the New Keynesian model with expectations based interest rate
rules under constant gain learning allows a comparison of policy outcomes
for varying degrees of commitment. \ The primary goals are to determine
whether there are gains to commitment and to find the optimal magnitude of
response to lagged output for the interest rate, which is parameterized by $%
\kappa $ in the policy rule (\ref{policy rule}). \ Since the full commitment
value of $\kappa =1.0$ is close to the bound for non-explosive and
determinate solutions for some reasonable parameter values, it is
questionable whether this policy remains optimal under learning. \ The
simulation results show that the modified commitment setting $\kappa =\beta $
minimizes the loss under least squares learning with constant gain. \
However, this and the following section show that consideration of parameter
uncertainty and the introduction of measurement error in the policymaker's
observations of public expectations complicates these results.

The choice of parameter values for this section follows McCallum and Nelson
(2004), MN04 in (\ref{table}), who study full commitment versus
discretionary outcomes for a number of different policy rules under full
rationality. \ They report gains from commitment over discretion for the
policymaker for a number of policy rules under different degrees of output
stabilization. \ The parameters are calibrated so that the time periods
correspond to quarterly data. \ The parameter values that determine the
behavior of the stochastic variables $g_{t}$ and $u_{t}$ in (\ref{IS}) and (%
\ref{Phillip's}) are as follows.%
\begin{equation*}
\begin{tabular}{|l|l|l|l|l|}
\hline
$\mu $ & $\rho $ & $\sigma _{w}$ & $\sigma _{g}$ & $\sigma _{u}$ \\ \hline
$0.95$ & $0.8$ & $0.005$ & $0.02$ & $0.005$ \\ \hline
\end{tabular}%
\end{equation*}%
Further, let $\alpha $ in the loss function (\ref{Loss}) be $\alpha =0.25,$
which is an intermediate value of those studied by McCallum and Nelson
(2004) and indicates equal weight on output and inflation stabilization. \ I
report results for values of the gain parameters from $\bar{\tau}$ $=0.01$
to $\bar{\tau}$ $=0.15$, which correspond to agents using a rolling window
of 200 and $3.\bar{3}$ years of past data, respectively. \ The smaller value
0.01 is the lowest considered by Orphanides and Williams (2006) while $\bar{%
\tau}=0.15$ is much larger than the values they consider, but in line with
the values used in other studies\footnote{%
See Evans and Honkapohja (1993) for an example where the gain parameter with
the optimal performance was 0.15. \ Orphanides and Williams (2002) use
semi-annual data so a gain of 0.05 is equivalent to a gain of 0.025 here.}.

\begin{center}
Figures 1 and 2 here
\end{center}

Figures 1 and 2 give sample paths for the output and inflation gaps for
simulations under recursive least squares with $\bar{\tau}=0.1$. \ Agents
begin the simulations with beliefs about the parameter values corresponding
to the REE, but with constant gain their estimates may continue to
fluctuate. \ For this simulation in Figure 1, the parameter $\kappa $ is set
to $\kappa =1.3$, which is in the range where the REE is determinate and
non-explosive. \ The RE value for $b_{x}\approx 0.98$ is close to one,
leading to the persistence in the series in Figure 1. \ In contrast, for
Figure 2, $\kappa =0.9$, the corresponding value is $b_{x}\approx 0.85,$ and
the policymaker is able to keep output and inflation closer to their target
values.

\begin{center}
Figure 3 here
\end{center}

To gain intuition about the problem facing the policymaker, Figure 3 graphs
the losses for $\kappa \in \lbrack 0.0,1.1]$ and $\alpha \in \lbrack
0.1,1.0] $ taking an average of the losses from 10,000 runs\footnote{%
For all reported results, the loss for each run is calculated over 200
periods using (\ref{Loss}) after 600 periods for initialization. \ The
initial values $\xi _{0}$ are set to their RE values. \ Since losses are
computed with discounting, longer runs do not provide much extra information.%
}.\ \ A number of observations can be made from Figure 3. \ First, the loss
is minimized in the neighborhood of $\kappa =1.0$ for all levels of alpha,
though there is an asymmetric change in the loss for $\kappa $ above and
below the minimizing value. \ Here is evidence in favor of partial
commitment if there is uncertainty about the true parameter values, since
accidentally using a policy equivalent to an excessively large $\kappa $
could lead to a very bad outcome. \ It is notable, however, that the
asymmetry is less severe for lower levels of alpha, in line with Proposition
1, which shows that a lower alpha raises the bound on $\kappa $\ for
determinacy and non-explosiveness. \ Again, a conservative central banker
reduces the risks associated with an excessive response to lagged output. \
Preliminary simulations indicate that this asymmetry is present under
rational expectations as well, though I leave a full investigation of this
issue for future work.

To give some precision to these observations, the table in Figure 4 gives
numerical values of losses for different values of $\kappa $ including those
for discretion, partial commitment and full commitment, with $\alpha =0.25$.
\ For this section the reader should focus on the first column of losses. \
Using the parameters of MN04, as in Figure 3, the minimum loss is achieved
at the partial commitment setting of $\kappa =0.99$. \ Hence, modified
commitment is optimal under least squares learning.

\begin{center}
Figure 4 here
\end{center}

\section{Imperfect Information for the Policymaker}

To this point I have made the assumption that the policymaker has perfect
knowledge of public expectations. \ Evans and Honkapohja (2004) state that
the results concerning expectational stability hold with white noise
measurement error in the expectations in the policy rule (\ref{policy rule}%
). \ Here, I study whether the introduction of such errors would impact the
policy outcomes under learning. \ Simulation results reported in this
section show that such errors can have a definite effect. \ In the presence
of such uncertainty, the policymaker should not respond to lagged output as
strongly as the full or modified commitment settings indicate.

To examine this issue, I include the error terms $\varepsilon _{x,t}$ and $%
\varepsilon _{\pi ,t}$, both white noise with variances $\sigma _{x}$ and $%
\sigma _{\pi }$, to the policy rule (\ref{policy rule}) to form\footnote{%
The form of this rule is identical to that found in Evans and Honkapohja
(2003).}%
\begin{equation}
i_{t}=\delta _{L}x_{t-1}+\delta _{\pi }\left( E_{t}^{\ast }\pi 
_{t+1}+\varepsilon _{\pi ,t}\right) +\delta _{x}\left( E_{t}^{\ast
}x_{t+1}+\varepsilon _{x,t}\right) +\delta _{g}g_{t}+\delta _{u}u_{t}.
\label{policy rule err}
\end{equation}%
Substituting this equation into (\ref{IS}) shows that the model is identical
to (\ref{matrix model}), but adds the additional noise $-\varepsilon 
_{x,t}-\left( 1+\dfrac{\lambda \beta }{\alpha +\lambda ^{2}}\right)
\varepsilon _{\pi ,t}$ to the preference shock $\widetilde{g}_{t}$. \
Introducing measurement error has the effect of adding unobservable shocks
to the demand side of the model. \ Besides measurement error, the extra
noise could be interpreted as any unseen factors affecting the connection
between the policymaker's instrument $i_{t}$ and the endogenous variables,
decisions within financial institutions being one of many examples. \ Since
the source of the errors is unimportant to the simulation results, I make
the simplifying assumption that $\varepsilon _{\pi ,t}=0$ and run
simulations for different levels of $\sigma _{x}$, adjusting $\sigma _{g}$
so that the magnitude of the sum of the shocks does not change\footnote{%
Formally, $\sigma _{g}^{2}+\sigma _{x}^{2}=\left( 0.02\right) ^{2}.$}.

\begin{center}
Figure 5 here
\end{center}

Introduction of the policy rule errors has an obvious affect on the public's
forecasting ability under learning. \ Figure 5 shows agents' parameter
estimates over time for $b_{x},b_{\pi },c_{x}$ and $c_{\pi }$ and
demonstrates this conclusion clearly. \ In the case of small or zero errors
in the interest rate rules, estimates of these parameters appear to be
constant\footnote{%
For example, of the four paramters, $c_{x}$ always has the largest mean
squared deviation (MSD), and for $\sigma _{x}=0.001$ its MSD is 0.0005, but
for $\sigma _{x}=0.01$, as shown in Figure 3, the MSD is 300 times higher at
0.15.} near their RE values, given by (\ref{coeff sols}). \ However, for
larger $\sigma _{x}=0.01$, as in Figure 5, the variation of these parameters
is quite visible. \ Greater variation in the parameter values leads to
higher volatility in expectations, see (\ref{exp}), and in output and
inflation, as well. \ Larger errors in the interest rate rule (\ref{policy
rule err}) lead to policy mistakes, making the task of estimating the values 
$\xi _{t}$ of the perceived model (\ref{PLM}) more difficult. \ The problem
is increasingly severe at higher levels of commitment, when the policymaker
is influencing public expectations, and for higher values of the gain
parameter $\bar{\tau}$, when agents are putting a great deal of emphasis on
recent information when estimating $\xi _{t}$.

In the cases with significant variation in $\xi _{t},$ the estimated values
in the PLM (\ref{PLM}) can easily slide into regions when the system appears
to be explosive. \ To ensure that the policy rules are operational, I impose
a bound or projection facility\footnote{%
This term originally comes from Marcet and Sargent's (1989) study of least
squares learning.} on the elements of $\xi _{t}$. \ Each element of $\xi 
_{t} $ is bounded by a unit interval centered at its RE value. \ The size of
the interval is set to be large enough so that the variation in the elements
of $\xi _{t}$ can have an impact on the endogenous variables but also to be
strict enough so that the simulated data does not exhibit unreasonably wild
swings. \ In Figure 6, the estimates of $c_{x}$ hit the bound a number of
times. \ Grandmont (1998) has questioned whether imposing a projection
facility excessively restricts agents' behavior, and our results are
sensitive to changes in the bound. \ Imposing a projection facility implies
that agents have a certain amount of faith that the economy will not
fluctuate excessively. \ Investigating whether and how such bounds should be
imposed and interpreted is a large area for future research. \ Here, I
simply note that expanding or removing the projection facility makes the
errors in the policy rule (\ref{policy rule err}) more of a concern for the
policymaker.

For each of the parameterizations in (\ref{table}), the table in Figure 4
reports losses for varying $\kappa $ for the case of no policy errors and
for errors with $\sigma _{x}=0.01$ where their impact is equal to that of
the preference shocks. \ For all three parameterizations under learning, the
modified commitment setting $\kappa =\beta =0.99$ remains optimal within the
class of rules (\ref{policy rule}), but within the class of rules (\ref%
{policy rule err}) with errors in the policy rule, a partial commitment
setting of $\kappa <\beta $ is optimal, $\kappa =0.94$ in the case of MN04.

The table in Figure 4 also reports losses in the discretionary case where $%
\kappa =0.0$ to ascertain the magnitude of the gains to commitment. \
Assuming rational expectations, McCallum and Nelson (2004) report a ratio of
losses between the discretionary and fully rational cases of 1.29. \ Under
learning, the ratios between discretion and both full and modified
commitment are larger at approximately 1.78. \ Using the alternative
parameterizations, CGG00 and W99, the qualitative conclusions are unchanged.

\begin{center}
Figure 6a and 6b here
\end{center}

The tables in Figures 6a and 6b report the loss minimizing $\kappa $'s under
learning for constant gain $\bar{\tau}=0.025$ and $\bar{\tau}=0.15$
respectively, for varying levels\footnote{%
The range for $\alpha $ corresponds to McCallum and Nelson (2004).} of the
policymaker preference parameter $\alpha $ and magnitude of policy rule
shocks $\sigma _{x}$. \ The surface of losses in Figure 3 shows that the
loss minimizing setting for $\kappa $ does not depend on $\alpha $, and,
indeed, modified commitment where $\kappa =0.99$ is best for all choices of $%
\alpha $ in the cases where there are no or minimal errors ($\sigma 
_{x}=0.001$) in the policy rule, according to the table in Figure 6a. \
However, for larger magnitudes of policy rule errors partial commitment
values of $\kappa <0.99$ minimize the policymaker's loss and vary inversely
with $\alpha $. \ The more significant the measurement error in the policy
rule and the greater the policymaker's emphasis on stabilizing output, the
less the interest rate rule should respond to lagged inflation. \ Results
using constant gain of $\bar{\tau}=0.01$ and $\bar{\tau}=0.05$ were nearly
identical to Figure 6a.

The case for partial commitment is even stronger for the higher constant
gain parameter $\bar{\tau}=0.15$ as in the table in Figure 6a. \ Modified
commitment is still best for any $\alpha $ in the absence of policy rule
errors, but partial commitment with $\kappa <0.99$ is optimal for minimal
errors ($\sigma _{x}=0.001$) if the the policymaker is sufficiently
concerned about output stabilization. \ With one exception\footnote{%
At $\alpha =0.01$ and $\sigma _{x}=0.01$ the loss minimizing $\kappa $ is
0.98 for $\bar{\tau}=0.025$ but 0.99 for $\bar{\tau}=0.15$. \ This result
was verified for 20,000 runs, but the difference in loss between $\kappa 
=0.99$ and $\kappa =0.98$ for the lower gain was less that $10^{-6}.$},
higher gain leads to a lower optimal level of commitment. \ The finding of a
negative $\kappa $ that minimizes the loss in the bottom right hand corner
is counter-intuitive, but the loss surface in the neighborhood of $\kappa =0$
is quite flat, and this outcome should be interpreted as a recommendation in
favor of discretionary policy for the indicated parameter values.

A significant deterioration of the performance of modified commitment
requires some combination of policymaker concern about output stabilization,
policy rule errors and/or high gain. \ Introduction of any one of these
factors by themselves does not alter greatly the conclusion that modified
commitment is preferable. \ For example, if policy rule errors are at their
maximum $\sigma _{x}=0.02$ while the policymaker is primarily focused on
inflation stabilization $\alpha =0.01$ and the gain is small $\bar{\tau}%
=0.025$, the loss minimizing $\kappa $ is 0.98 (Figure 6a) and the gain over
modified commitment $\kappa =0.99$ is less than $10^{-5}$. \ However,
increasing $\alpha $ with the higher policy rule errors changes this
conclusion making partial commitment better and raising the gain parameter
in addition lowers the loss minimizing $\kappa $ even more.

\begin{center}
Figure 7 here
\end{center}

To clarify the results concerning the optimality of partial commitment for
this class of interest rate rules (\ref{policy rule err}), I closely examine
four cases under modified commitment where the standard deviation of the
policy rule errors takes the value of $\sigma _{x}=0.001$ or 0.01, and the
gain parameter takes the values $\bar{\tau}=0.025$ and $0.15$.\ \ The table
in Figure 7 reports the mean deviation of the estimates of $b_{x},b_{\pi 
},c_{x}$ and $c_{\pi }$ from their RE values, their squared deviation from
the mean of the simulated values, the mean squared error of the forecasts
for output and inflation, and the skewness of the losses over the 10,000
runs. \ Orphanides and Williams (2002) report systematic bias in agents'
estimates even in the absence of stochastic shocks in the model. \ Here,
bias and fluctation in the estimates of the parameters are evident and
coincide with larger forecast errors.

The first two columns are cases with minimal policy rule errors when
modified commitment $\kappa =\beta =0.99$ is best. \ In both cases, the bias
and fluctation in the parameter estimates and the forecast errors are
smaller than the cases in the third and fourth columns. \ The presence of
significant policy rule errors has a detrimental effect on agents
forecasting ability, most notably for output. \ The losses in the first two
columns show positive skewness, as expected given the asymmetry in Figure 3,
though it is not sufficient to affect the optimality of modified commitment.
\ The most notable entry in Figure 7 is the high, positive skewness in the
case with large policy rule errors and gain, $\left( \sigma _{x},\bar{\tau}%
\right) =\left( 0.01,0.15\right) $. \ In this case, the prevalence of some
losses far above the mean makes a setting of $\kappa =0.91$ (see Figure 6b)
better than modified commitment. \ However, comparing the second and third
columns, the case where $\left( \sigma _{x},\bar{\tau}\right) =\left(
0.001,0.15\right) $ corresponds to a loss minimizing $\kappa $ of $\kappa 
=0.99$ while the case where $\left( \sigma _{x},\bar{\tau}\right) =\left(
0.01,0.025\right) $ has a lower loss minimizing $\kappa =0.94$ even though
the skewness is less than in the second column. \ Hence, larger bias and/or
variation in the parameter estimates or the larger policy rule errors in the
third column must also be a factor behind the optimal $\kappa $ being below
modified commitment.

Simulations of the model under partial commitment in the case where $\left(
\sigma _{x},\bar{\tau}\right) =\left( 0.01,0.025\right) $ for a fixed $\xi 
$%
, further disentangle the influence of these factors. \ Fixing the agents
estimate of $\xi $ at the biased value from Figure 7 yielded a loss of
0.09120, which is a slight improvement over the 0.0925 shown in Figure 4 but
not nearly as good as the loss achieved at the optimal $\kappa =0.94$. \
Fixing the estimate of $\xi $ at the rational expectations value yielded a
loss of 0.09219, which is worse than the outcome under the biased $\xi $, so
the bias actually improved the policy outcome in the presence of policy rule
errors. \ Hence, the fluctuations in the estimates of $\xi $ play a role in
the deterioration of the performance of modified commitment, but the
presence of the policy rule errors are a major factor. \ 

\begin{center}
Figure 8 here
\end{center}

The deterioration of policy outcomes at modified commitment for larger
policy rules errors and higher gain is evident from the distribution of the
losses shown in Figure 8. \ These curves track the height of the midpoints
of a histogram over 100 intervals of width 0.003. \ For the case where $%
\left( \sigma _{x},\bar{\tau}\right) =\left( 0.01,0.15\right) $ the mean is
to the right, and the tail is fatter compared to the other distributions. \
The jump at 0.3 for this case indicates there are a number of cases with
losses above the maximum of 0.3 as well..

Close examination of these four cases shows that larger policy rule errors,
greater variation in the estimates of $\xi $ and higher gain lead to
increased incidence of large losses and worse policy outcome under modified
commitment, necessitating a partial commitment setting for $\kappa <\beta $.
\ Conversely, for smaller policy rule shocks and low gain, the fluctuations
in $\xi $ are minimal, the asymmetry of the losses seen in Figure 3 do not
play a significant role and the modified commitment setting $\kappa =\beta $
is optimal under learning.

The interpretation of the results with respect to $\sigma _{x}$ is
straightforward. \ Larger errors in the policy rule indicate a lower level
of knowledge on the part of the policymaker and a public who is less able to
correctly learn the values of the forecasting model. \ Policymakers cannot
set interest rates as accurately and expectations show more instability
leading to worse outcomes from a policy perspective.

The impact of a higher $\bar{\tau}$ is more open to interpretation. \ The
gain parameter can represent the lags of data used by agents in their
estimation, the rationality of agents and/or the credibility of the
policymaker. \ A marked policy deterioration only appears when $\bar{\tau}$
is raised to 0.15, a relatively high value indicating that agents are using
only 3.3 years of data to make their estimates. \ This higher gain could be
reasonable if there is a perceived shift of regime in the economy or if the
policymaker has very low credibility. \ The policymaker is less able to
commit when he or she has low credibility or poor information about the
connection between the policy instrument and the other endogenous variables
in the economy.

The conclusion that a policymaker with weak credibility should not
aggressively commit rests on the assumption of constant gain, however. \ If
the model is expanded to allow the gain parameter to fall as the policymaker
gains credibility, then the policymaker may have incentive to stick to a
higher level of commitment. \ Such a model is studied in Waters (2007) who
uses an endogenous gain mechanism introduced by Marcet and Nicolini (2004).

\section{Conclusion}

Relaxing the assumption of full rationality for the formation of
expectations raises a number of related issues for the design of monetary
policy rules. \ One concern is whether a rule is stable under learning while
another is whether a rule minimizes the policymaker's loss under learning. \
Evans and Honkapohja (2003, 2006) introduce a class of expectations based
interest rate rules and study the former issue. \ The present work extends
their analysis to a broader class of rules allowing for varying degrees of
response to lagged variables and goes on to study the latter question for
such rules.

Within the present framework, one can meaningfully discuss varying levels of
commitment. \ The greater the policymaker's degree of commitment, the more
he or she will adjust the interest rate in response to lagged output to
affect public expectations. \ The commitment optimum under rational
expectations\ or full commitment is a special case as is the modified
commitment value advocated by Blake (2001). I provide a condition for
non-explosive and determinate equilibria for this class of rules and show
that expectational stability holds for any non-negative response to lagged
output.

Simulation results show that modified commitment minimizes the policymaker's
loss under least squares learning in the baseline case. \ However, parameter
uncertainty, errors in the policy rule and high gain are all factors that
could, in combination, make a policy of partial commitment best implying a
lower response to lagged output than modified commitment. \ An excessive
response to lagged output can lead to very bad outcomes for the policymaker,
so he or she should act to minimize that possibility. \ The presence of a
conservative central banker makes such issues less problematic, however.

The gain parameter has important economic interpretations in terms of the
lags used to form expectations, the rationality of the agents and the
credibility of the policymaker. \ When agents are using only recent data
(high gain), indicating a low credibility, the case for partial commitment
becomes stronger.

Expectations based interest rate rules continue to demonstrate desirable
characteristics for monetary policy, but the present work does offer a few
caveats for practical implementation. \ For an extended range of responses
to lagged output they exhibit non-explosive, determinate and expectationally
stable solutions. \ Some degree of commitment minimizes the loss for the
policymaker when agents use a learning mechanism to form expectations. \
However, it is often the case that the policymaker should be cautious about
the commitment optimum under rational expectations and engage in partial
commitment. \ The prerequisites for partial commitment, parameter
uncertainty and policy rule errors, clearly exist to some degree, so the
policymaker must be vigilant.

These results point to some important areas for future work, namely
determining the extent to which parameter uncertainty and policy rule errors
are present in estimated models of monetary policy. Furthermore, credibility
is an important aspect of monetary policy but has proved to be an elusive
feature to model formally. \ Interpreting the gain parameter as an
indication of credibility may open new avenues for research on monetary
policy. \ The present paper gives practical guidance on the design of rules
and areas that need attention for deeper understanding of monetary policy.

\bigskip

{\LARGE Appendix A}

\textbf{Proof of Proposition 1:}

\begin{proof}
To demonstrate determinacy, first, define a new variable $x_{t}^{L}$ for the
lag of the output gap such that $x_{t+1}^{L}=x_{t}$. \ Under rational
expectations $E_{t}\pi _{t+1}=\pi _{t+1}+\varepsilon _{t+1}$ for some 
\textit{iid}, mean 0 shock $\varepsilon _{t}$. \ The the model (\ref{matrix
model}) can be written in the form 
\begin{equation*}
\left( 
\begin{array}{ccc}
1 & 0 & \dfrac{-\alpha \kappa }{\alpha +\lambda ^{2}} \\ 
0 & 1 & \dfrac{-\lambda \alpha \kappa }{\alpha +\lambda ^{2}} \\ 
1 & 0 & 0%
\end{array}%
\right) \left( 
\begin{array}{c}
x_{t} \\ 
\pi _{t} \\ 
x_{t}^{L}%
\end{array}%
\right) =\left( 
\begin{array}{ccc}
0 & \dfrac{-\lambda \beta }{\alpha +\lambda ^{2}} & 0 \\ 
0 & \dfrac{\alpha \beta }{\alpha +\lambda ^{2}} & 0 \\ 
0 & 0 & 1%
\end{array}%
\right) \left( 
\begin{array}{c}
x_{t+1} \\ 
\pi _{t+1} \\ 
x_{t+1}^{L}%
\end{array}%
\right) +K\left( 
\begin{array}{c}
u_{t} \\ 
g_{t} \\ 
\varepsilon _{t}%
\end{array}%
\right) .
\end{equation*}%
Multiplying both sides by the inverse of the 3x3 matrix on the left hand
side gives the model the following form.%
\begin{equation*}
\left( 
\begin{array}{c}
x_{t} \\ 
\pi _{t} \\ 
x_{t}^{L}%
\end{array}%
\right) =J\left( 
\begin{array}{c}
x_{t+1} \\ 
\pi _{t+1} \\ 
x_{t+1}^{L}%
\end{array}%
\right) +\tilde{K}\left( 
\begin{array}{c}
u_{t} \\ 
g_{t} \\ 
\varepsilon _{t}%
\end{array}%
\right)
\end{equation*}%
where $\tilde{K}$ is the appropriate transformation of $K$. \ The condition
for determinacy (see Blanchard and Kahn (1980) and Evans and Honkapohja
(2006)) is that the matrix $J$ must have two eigenvalues with modulus less
than one and one eigenvalue with modulus greater than one.

A straightforward computation shows that one eigenvalue $e$\ of $J$ is 0,
while the other two must satisfy%
\begin{equation*}
e^{2}-\left( \beta +\frac{\alpha +\lambda ^{2}}{\alpha \kappa }\right) e+%
\frac{\beta }{\kappa }=0.
\end{equation*}%
\ \ For one root to be greater than one and one root to be less than one the
right hand side of the equation above must be negative for $e=1$, which is
equivalent to the condition $\kappa <1+\dfrac{\lambda ^{2}}{\alpha \left(
1-\beta \right) }$.

Non-explosiveness requires that the negative solution of $b_{x}$ in the
quadratic condition (\ref{quad cond}) be less than one. \ Since the positive
root is always greater than one, the non explosiveness condition requires
that the left hand side of (\ref{quad cond}) be less than zero at $b_{x}=1$
which is true for $\kappa <1+\dfrac{\lambda ^{2}}{\alpha \left( 1-\beta 
\right) }$ as well.
\end{proof}

\bigskip

{\LARGE Appendix B}

\textbf{Proof of Proposition 2:}

To prove proposition 2, I first prove the following lemma.

\begin{lemma}
Given $\kappa \geq 0$ and the associated REE value of $b_{x}$,

i) $b_{x}$ is real,

ii) $b_{x}-\kappa \leq 0$,

iii) $b_{x}-\kappa >-1.$
\end{lemma}

\begin{proof}
The relevant root of (\ref{quad cond}) is%
\begin{equation*}
b_{x}=\left( 2\beta \right) ^{-1}\left( \frac{\lambda ^{2}+\alpha }{\alpha }%
+\beta \kappa -\sqrt{\left( \frac{\lambda ^{2}+\alpha }{\alpha }+\beta 
\kappa \right) ^{2}-4\beta \kappa }\right) ,
\end{equation*}%
which is real if the radicand $rad$ is positive. \ The radicand can be
written%
\begin{equation*}
rad=\beta ^{2}\kappa +2\beta \kappa \left( \dfrac{\lambda ^{2}-\alpha }{%
\alpha }\right) +\left( \frac{\lambda ^{2}+\alpha }{\alpha }\right) ^{2},
\end{equation*}%
and is minimized at $\kappa =-\dfrac{1}{\beta }\left( \dfrac{\lambda 
^{2}-\alpha }{\alpha }\right) $ so, substituting, the minimum value attained
by $rad$ is $\dfrac{4\lambda ^{2}}{\alpha }>0.$ \ Hence, the radicand is
always positive for positive parameters, and $b_{x}$ is real.

The difference $b_{x}-\kappa $ can be written%
\begin{equation*}
b_{x}-\kappa =\left( 2\beta \right) ^{-1}\left( \frac{\lambda ^{2}+\alpha }{%
\alpha }-\beta \kappa -\sqrt{\left( \frac{\lambda ^{2}+\alpha }{\alpha }%
-\beta \kappa \right) ^{2}-4\beta \kappa }\right) .
\end{equation*}%
Therefore $b_{x}-\kappa \leq 0$ if $4\beta \kappa >0$ which is true if $%
\kappa \geq 0$.

The condition in \textit{iii)}, $b_{x}-\kappa >-1$, is equivalent to 
\begin{equation*}
2\beta +\frac{\lambda ^{2}+\alpha }{\alpha }-\beta \kappa >\sqrt{\left( 
\frac{\lambda ^{2}+\alpha }{\alpha }-\beta \kappa \right) ^{2}-4\beta \kappa 
}.
\end{equation*}%
Squaring both sides and simplifying yields another equivalent condition%
\begin{equation*}
\beta +\frac{\lambda ^{2}+\alpha }{\alpha }>-\kappa \left( 1-\beta \right) ,
\end{equation*}%
which is true for $\kappa \geq 0$, $\beta <1$ and positive parameters.
\end{proof}

To study expectational stability of (\ref{matrix model}), consider the
general model%
\begin{equation*}
y_{t}=ME_{t}^{\ast }y_{t+1}+Ny_{t-1}+Pv_{t}
\end{equation*}%
with $v_{t}=Fv_{t-1}$ from (\ref{shocks}). \ Expectations $E_{t}^{\ast 
}y_{t+1}$ are formed according to (\ref{exp}) and substituting them into
equation above gives the ALM in the form $y_{t}=T\left(
a_{t-1},b_{t-1},c_{t-1}\right) z_{t-1}$ for the map 
\begin{equation*}
T\left( a,b,c\right) =\left( M\left( a+ba\right) ,\text{ }Mb^{2}+N,\text{ }%
M\left( bc+F\right) +P\right)
\end{equation*}%
suppressing the time subscripts in the $T-$ map. \ For the model (\ref%
{matrix model}) with REE (\ref{RE sols}) the matrices are such that $%
M=\left( 
\begin{array}{cc}
0 & \dfrac{-\lambda \beta }{\alpha +\lambda ^{2}} \\ 
0 & \dfrac{\alpha \beta }{\alpha +\lambda ^{2}}%
\end{array}%
\right) $ and $B=\left( 
\begin{array}{cc}
b_{x} & 0 \\ 
b_{\pi } & 0%
\end{array}%
\right) .$ \ Expectational stability is defined according to the matrix
differential equation%
\begin{equation*}
\frac{d}{ds}\left( a,b,c\right) =T\left( a,b,c\right) -\left( a,b,c\right)
\end{equation*}%
where $s$ is notional time, distinct from the periods denoted by $t$. \ The
REE solutions are fixed points of the differential equation where $T\left(
a,b,c\right) =\left( a,b,c\right) ,$ and expectational stability is defined
in terms of stability of the differential equation.

The associated REE from (\ref{RE sols}) has the form%
\begin{equation*}
y_{t}=By_{t-1}+Cv_{t}.
\end{equation*}%
The conditions for expectational stability (see Evans and Honkapohja (2001),
chapter 10) are that the matrices%
\begin{eqnarray*}
DT_{a} &=&M\left( I+B\right) \\
DT_{b} &=&B^{\prime }\otimes M+I\otimes MB \\
DT_{c} &=&F^{\prime }\otimes M+I\otimes MB
\end{eqnarray*}%
must have eigenvalues with modulus less than one. \ 

The proof of proposition 2 follows.

\begin{proof}
The matrices for the condition for expectational stability are%
\begin{eqnarray*}
DT_{a} &=&%
\begin{pmatrix}
\dfrac{-b_{\pi }\lambda \beta }{\alpha +\lambda ^{2}} & \dfrac{-\lambda 
\beta }{\alpha +\lambda ^{2}} \\ 
\dfrac{b_{\pi }\alpha \beta }{\alpha +\lambda ^{2}} & \dfrac{\alpha \beta }{%
\alpha +\lambda ^{2}}%
\end{pmatrix}
\\
&& \\
DT_{b} &=&\left( 
\begin{array}{cccc}
\dfrac{-b_{\pi }\lambda \beta }{\alpha +\lambda ^{2}} & \dfrac{-b_{x}\lambda
\beta }{\alpha +\lambda ^{2}} & 0 & \dfrac{-b_{\pi }\lambda \beta }{\alpha 
+\lambda ^{2}} \\ 
\dfrac{b_{\pi }\alpha \beta }{\alpha +\lambda ^{2}} & \dfrac{b_{x}\alpha 
\beta }{\alpha +\lambda ^{2}} & 0 & \dfrac{b_{x}\alpha \beta }{\alpha 
+\lambda ^{2}} \\ 
0 & 0 & \dfrac{-b_{\pi }\lambda \beta }{\alpha +\lambda ^{2}} & 0 \\ 
0 & 0 & \dfrac{b_{\pi }\alpha \beta }{\alpha +\lambda ^{2}} & 0%
\end{array}%
\right) \\
&& \\
DT_{c} &=&\left( 
\begin{array}{cccc}
\dfrac{-b_{\pi }\lambda \beta }{\alpha +\lambda ^{2}} & \dfrac{-\mu \lambda
\beta }{\alpha +\lambda ^{2}} & 0 & 0 \\ 
\dfrac{b_{\pi }\alpha \beta }{\alpha +\lambda ^{2}} & \dfrac{\mu \alpha 
\beta }{\alpha +\lambda ^{2}} & 0 & 0 \\ 
0 & 0 & \dfrac{-b_{\pi }\lambda \beta }{\alpha +\lambda ^{2}} & \dfrac{-\rho 
\lambda \beta }{\alpha +\lambda ^{2}} \\ 
0 & 0 & \dfrac{b_{\pi }\alpha \beta }{\alpha +\lambda ^{2}} & \dfrac{\rho 
\alpha \beta }{\alpha +\lambda ^{2}}%
\end{array}%
\right) .
\end{eqnarray*}%
The eigenvalues of $DT_{a}$ are 0 and $\dfrac{\alpha \beta }{\alpha +\lambda 
^{2}}\left( 1-\dfrac{\lambda }{\alpha }b_{\pi }\right) $. \ Using the second
equation from (\ref{coeff sols}), the non-zero eigenvalue is $\dfrac{\alpha 
\beta }{\alpha +\lambda ^{2}}\left( 1+b_{x}-\kappa \right) $, which is less
than one since $\dfrac{\alpha }{\alpha +\lambda ^{2}},\beta <1$ and $%
b_{x}-\kappa \leq 0$ from Lemma 3. \ 

The the matrix $DT_{b}$ has two eigenvalues equal to 0. \ It also has
eigenvalues $\dfrac{\alpha \beta }{\alpha +\lambda ^{2}}\left( b_{x}-\dfrac{%
\lambda }{\alpha }b_{\pi }\right) $ and $\dfrac{-b_{\pi }\lambda \beta }{%
\alpha +\lambda ^{2}}$. \ Using the second equation from (\ref{coeff sols}),
the two non-zero eigenvalues are $\dfrac{\alpha \beta }{\alpha +\lambda ^{2}}%
\left( b_{x}+b_{x}-\kappa \right) $, and $\dfrac{\alpha \beta }{\alpha 
+\lambda ^{2}}\left( b_{x}-\kappa \right) ,$ the second of which is less
than one, since $\dfrac{\alpha }{\alpha +\lambda ^{2}},\beta <1$ and $%
b_{x}-\kappa \leq 0$ from Lemma 3. \ The condition on the first non-zero
eigenvalue of $DT_{b}$, $\dfrac{\alpha \beta }{\alpha +\lambda ^{2}}\left(
b_{x}+b_{x}-\kappa \right) <1$ is equivalent to 
\begin{equation*}
1-\frac{1}{2}\left( \dfrac{\alpha }{\alpha +\lambda ^{2}}\right) \left[ 
\sqrt{\left( \frac{\lambda ^{2}+\alpha }{\alpha }+\beta \kappa \right) 
^{2}-4\beta \kappa }+\sqrt{\left( \frac{\lambda ^{2}+\alpha }{\alpha }-\beta
\kappa \right) ^{2}-4\beta \kappa }\right] <1,
\end{equation*}%
using the expressions for $b_{x}$ and $b_{x}-\kappa $ in Lemma 3. \ The
condition is satisfied as long as both radicands from the roots of $b_{x}$
and $b_{x}-\kappa $ are positive. \ Lemma 3 proves that $b_{x}$ is real,
hence $b_{x}-\kappa $ is as well, both radicands are positive, and the
eigenvalue is less than one.

The matrix $DT_{c}$ has two repeated eigenvalues of 0 and $\dfrac{\alpha 
\beta }{\alpha +\lambda ^{2}}\left( \mu -\dfrac{\lambda }{\alpha }b_{\pi 
}\right) =\dfrac{\alpha \beta }{\alpha +\lambda ^{2}}\left( \mu 
+b_{x}-\kappa \right) $, using the second equation from (\ref{coeff sols})
again. \ The non-zero eigenvalue is less that one since $\dfrac{\alpha }{%
\alpha +\lambda ^{2}},\beta ,\mu <1$ and $b_{x}-\kappa <0$ from Lemma 3.

The lemma also guarantees the $b_{x}-\kappa >-1$ so all the non-zero
eigenvalues for the three matrices are greater than $-1$ as well.

Therefore, the moduli of all the eigenvalues of $DT_{a},DT_{b}$ and $DT_{c}$
are all less than one and the model (\ref{matrix model}) is expectationally
stable for any REE.
\end{proof}

\bigskip

{\Large References}

\ \newline
Barro, Robert J. and David B. Gordon (1983) A Positive Theory of Monetary
Policy in a Natural Rate Model. \textit{Journal of Political Economy} 91,
589-610.\bigskip \newline
Beechey, Meredith (2004) Excess Sensitivity and Volatility of Long Interest
Rates: The Role of Limited Information in Bond Markets. Working Paper,
UC-Berkeley.\bigskip \newline
Blake, A. P. (2001) A `Timeless Perspective' on Optimality in
Forward-Looking Rational Expectations Models. Working Paper, National
Institute for Economic and Social Research.\bigskip \newline
Bullard, James and In-Koo Cho (2005) Escapist Policy Rules. \ \textit{%
Journal of Economic Dynamics and Control}, 29(11), 1841-65.\bigskip \newline
Bullard, James and Kaushik Mitra (2002) Learning about Monetary Policy
Rules. \textit{Journal of Monetary Economics} 49, 1105-1129.\bigskip \newline
Carceles-Poveda, Eva and Chryssi Giannitsarou (2006) Asset Pricing with
Adaptive Learning. Working Paper.\bigskip \newline
Clarida, Richard and Jordi Gali and Mark Gertler (1999) The Science of
Monetary Policy: \ A New Keynesian Perspective. \textit{Journal of Economic
Literature} 37, 1661-1707.\bigskip \newline
Clarida, Richard and Jordi Gali and Mark Gertler (2000) Monetary Policy
Rules and Macroeconomic Stability: \ Evidence and Some Theory. \textit{%
Quarterly Journal of Economics} 115, 147-180.\bigskip \newline
Evans, George and Seppo Honkapohja (1993) Adaptive Forecasts, Hysteresis and
Endogenous Fluctuations. \textit{Federal Reserve Bank of San Francisco
Economic Review} 1993(1), 3-13.\bigskip \newline
Evans, George and Seppo Honkapohja (2001)\ \textit{Learning and Expectations
in Macroeconomics}. Princeton University Press, Princeton, NJ.\bigskip 
\newline
Evans, George and Seppo Honkapohja (2003) Expectations and the Stability
Problem for Optimal Monetary Policies. \textit{Review of Economic Studies}
70, 807-824.\bigskip \newline
Evans, George and Seppo Honkapohja (2006) Monetary Policy, Expectations and
Commitment. \textit{Scandinavian Journal of Economics} 108, 15-38.\bigskip 
\newline
Evans, George and Bruce McGough (2004) Monetary Policy, Indeterminacy and
Learning. \textit{Journal of Economic Dynamics and Control }29,
1809-1840.\bigskip \newline
Evans, George and Bruce McGough (2006) Optimal Constrained Interest Rate
Rules. \textit{Journal of Money, Credit and Banking, }forthcoming\textit{.}%
\bigskip \newline
Fuhrer, Jeffrey and Mark Hooker (1993) Learning about Monetary Regime Shifts
in an Overlapping Wage Contract Model. \textit{Journal of Economic Dynamics
and Control} 17(4), 531-53.\bigskip \newline
Honkapohja, Seppo and Kaushik Mitra (2003) Learning with Bounded Memory in
Stochastic Models. \textit{Journal of Economic Dynamics and Control} 27,
1437-1457.\bigskip \newline
Honkapohja, Seppo and Kaushik Mitra (2004) Are Non-Fundamental Equilibria
Learnable in Models of Monetary Policy? \textit{Journal of Monetary Economics%
} 51(8), 1743-1770.\bigskip \newline
Howitt, Peter (1992) Interest Rate Control and Nonconvergence to Rational
Expectations. \textit{Journal of Political Economy} 100, 776-800.\bigskip 
\newline
Jensen, Christian and Bennet McCallum (2002) The Non-Optimality of Proposed
Monetary Policy Rules Under Timeless-Perspective Commitment. \textit{%
Economics Letters}, 77(2), 163-168.\bigskip \newline
Kydland, Finn E. and Edward C. Prescott (1977) Rules Rather than Discretion:
\ The Inconsistency of Optimal Plans. \textit{Journal of Political Economy}
85, 473-493.\bigskip \newline
Marcet, Albert and Juan Nicolini (2004) Recurrent Hyperinflations and
Learning. \textit{American Economic Review} 93(5), 1476-1498.\bigskip 
\newline
McCallum, Bennet T. (1983) On Non-uniqueness in Linear Rational Expectations
Models: \ An Attempt at Perspective. \textit{Journal of Monetary Economics}
11, 139-168.\bigskip \newline
McCallum, Bennet T. (1997) The Role of the Minimum State Variables Criterion
in Rational Expectations Models. \ \textit{Journal of International Tax and
Finance} 6(4), 621-639.\bigskip \newline
McCallum, Bennet T. and Edward Nelson (2004) Timeless Perspective vs.
Discretionary Monetary Policy in Forward Looking Models. \textit{Federal
Reserve Bank of St. Louis Review} 86(2), 43-56..\bigskip \newline
Orphanides, Athanios and John C. Williams (2002) Imperfect Knowledge,
Inflation Expectations and Monetary Policy. NBER Working Paper 9884.\bigskip 
\newline
Orphanides, Athanios and John C. Williams (2006) Monetary Policy with
Imperfect Knowledge. \textit{Journal of the European Economic Association }%
4(2-3), 366-375. \bigskip \newline
Evans, George and Seppo Honkapohja (2006) Monetary Policy, Expectations and
Commitment. \textit{Scandinavian Journal of Economics} 108, 15-38.\bigskip 
\newline
Walsh, Carl E. (2003) \textit{Monetary Theory and Policy}. MIT Press,
Cambridge, MA.\bigskip \newline
Waters, George (2007) Regime Changes, Learning and Monetary Policy. \textit{%
Journal of Macroeconomics} 29(2), 255-282.\bigskip \newline
Woodford, Michael (1999a) Optimal Monetary Policy Inertia. \textit{The
Manchester School, Supplement 67, 1-35.}\bigskip \newline
Woodford, Michael (1999b)Commentary: \ How Should Monetary Policy Be
Conducted in an Era of Price Stability? In \textit{New Challenges for
Monetary Policy}, pp. 277-316, Federal Reserve Bank of Kansas City.\bigskip 
\newline
Woodford, Michael (2003) \textit{Interest and Prices}. Princeton University
Press, Princeton, NJ.

\pagebreak

\end{doublespace}

\end{document}
